We begin by fixing the notation for counting processes,
reviewing the concepts from the general theory that the rest of the chapter uses.
Our main reference is \cite[Chapter 14]{DVJ08int}.

Let $\numEventTypes$ be a positive integer.
For each $e$ ranging from $1$ to $\numEventTypes$,
let $\arrivalTimes_{j}$, $j = 1, 2, \dots$, be a strictly increasing sequence of positive random times,
and assume that
\begin{equation}\label{eq.noCommonJumps}
 \arrivalTimes_{j} \neq \arrivalTimes[\eone]_{j\derivative}
 \qquad \text{whenever } (e,j) \neq (\eone, j\derivative).
\end{equation}
No two events ever occur at the same instant.
This is a modelling assumption and not a triviality:
it will be needed again for the second-order structure of Section \ref{sec.secondOrder}
and for the goodness-of-fit test of Section \ref{sec.simulation},
and it is the hypothesis that fails first when a model is fitted to recorded data,
where a single aggressive order is reported as several fills sharing a timestamp.

Under \eqref{eq.noCommonJumps},
\begin{equation*}
 \countingProc(t) := \sum_{j} \indicator{\arrivalTimes_{j} \leq t},
 \qquad t \geq 0,
\end{equation*}
is a non-decreasing right-continuous process;
we call $\countingProc$ the counting process associated with the sequence $(\arrivalTimes_{j})_j$.
Notice that $(\arrivalTimes_{j})_j$ can be retrieved from $\countingProc$ by
\begin{equation*}
 \arrivalTimes_{j} = \inf \left\lbrace t>0: \, \countingProc(t) \geq j \right\rbrace;
\end{equation*}
hence there is a one-to-one correspondence between $\countingProc$ and $(\arrivalTimes_{j})_j$.

For $t>0$ we define
\begin{equation*}
 \delta\countingProc(t) := \lim_{h\downarrow 0} \Big( \countingProc(t) - \countingProc(t-h) \Big),
\end{equation*}
and we notice that $\delta \countingProc(t) = 1$ if and only if $t = \arrivalTimes_{j}$ for some $j$,
and $\delta \countingProc(t) = 0$ otherwise.

The $\numEventTypes$-dimensional vector
$\multiCountingProc(t) = (\countingProc[1](t), \dots , \countingProc[\numEventTypes](t))$
is referred to as the multivariate counting process associated with the $\numEventTypes$ sequences
$(\arrivalTimes_{j})_j$, $e = 1, \dots, \numEventTypes$.
Let $\groundProc(t) := \countingProc[1](t) + \dots + \countingProc[\numEventTypes](t)$
be the ground process of $\multiCountingProc$,
and let
\begin{equation*}
 \arrivalTimes[ ]_{n} := \inf \left\lbrace t>0: \, \groundProc(t) \geq n\right\rbrace,
 \qquad n = 1, 2, \dots,
\end{equation*}
be the ordered sequence of random times stemming from the union
$\lbrace \arrivalTimes_{j}: \, j = 1,2,\dots; \,\, e = 1,\dots,\numEventTypes\rbrace$.
By defining, for $n = 1, 2, \dots$,
\begin{equation*}
 \event := \sum_{e=1}^{\numEventTypes} e \,
 \indicator{\delta \groundProc(\arrivalTimes[ ]_{n}) = \delta\countingProc(\arrivalTimes[ ]_{n})},
\end{equation*}
we have that the pair $(\arrivalTimes[ ]_{n}, \event[n])$ equivalently characterises the
multivariate counting process, because
\begin{equation}\label{eq.NTE_counting_proc}
 \countingProc(t) = \sum_{n} \indicator{\arrivalTimes[ ]_{n}\leq t, \, \event[n] = e},
\end{equation}
for all $t>0$ and all $e = 1,\dots, \numEventTypes$.

We can interpret this construction by saying that the index $e$ labels $\numEventTypes$ types
of events that occur in time,
and that $\countingProc(t)$ counts the number of events of type $e$ that have occurred by time $t$.
The two representations are the two we shall use throughout:
$\multiCountingProc$ when we compute,
and $(\arrivalTimes[ ]_{n}, \event[n])$ when we simulate,
since a simulator produces one time and one label at a time.

\begin{example}[``Poisson process'']
 Let $\tau_e^j$, $j = 1,2,\dots$, $e = 1,\dots, \numEventTypes$, be independent random variables
 such that $\tau_e^j$ is exponentially distributed with parameter $\intensity>0$.
 Let $\arrivalTimes_{j} := \sum_{k\leq j} \tau_e^{k}$,
 and notice that $\arrivalTimes_{j}$ has probability density function
 \begin{equation*}
  f_{e,j}(t) = \frac{\intensity^j}{(j-1)!}t^{j-1} e^{-\intensity t} \indicator{t>0}.
 \end{equation*}
 Then the multivariate counting process $\multiCountingProc$ associated with the arrival times
 $\arrivalTimes_{j}$ is called $\numEventTypes$-dimensional Poisson process of rates
 $\intensity[1], \dots, \intensity[\numEventTypes]$.
 This name is justified as follows.
 Since $\lbrace \countingProc(t)\geq j\rbrace = \lbrace \arrivalTimes_{j} \leq t\rbrace$,
 we have that $\frac{d}{dt} \Prob(\countingProc(t) \geq j) = f_{e,j} (t)$.
 On the other hand, if we define
 \begin{equation*}
  G_{e,j}(t) := \sum_{k\geq j} \frac{(\intensity t)^k}{k!}e^{-\intensity t},
 \end{equation*}
 we also have that $\frac{d}{dt} G_{e,j}(t) = f_{e,j}(t)$, by telescopic sum.
 Since $G_{e,j}(0) = \Prob (\countingProc(t)\geq 0)$,
 we deduce that $\Prob(\countingProc(t) \geq j) = G_{e,j}(t)$, and that
 \begin{equation*}
  \Prob\left( \countingProc(t) = j\right) =
  \frac{(\intensity t)^j}{j!} \exp\left( -\intensity t\right).
 \end{equation*}
 Therefore for every $t$, $\countingProc(t)$ is a Poisson random variable of parameter
 $\intensity t$,
 and the ground process $\groundProc$ of $\multiCountingProc$ is such that for every $t$,
 $\groundProc(t) \sim \text{Pois}(\intensity[1]t + \dots + \intensity[\numEventTypes]t)$.
\end{example}

The minimal filtration to which a multivariate counting process $\multiCountingProc$ is adapted
-- and which satisfies the usual conditions of completeness and right-continuity --
is called the internal history of $\multiCountingProc$.
Any other filtration to which $\multiCountingProc$ is adapted is called a history of
$\multiCountingProc$, and it must be a superset of the internal history.
The distinction will matter in Section \ref{sec.readingOrderFlow}:
an observer who sees less than the internal history has a different compensator for the same
process, and the quantity he can estimate is not the one the model specifies.

\begin{defi}\label{def.compensator}
 Let $\stochasticBase$ be a filtered probability space where the multivariate counting process
 $\multiCountingProc$ is defined,
 and assume that $\filtrationF$ is a history of $\multiCountingProc$.
 We say that the $\numEventTypes$-dimensional stochastic process
 $\compensator = (\compensator_1, \dots , \compensator_{\numEventTypes})$
 is an $\filtrationF$-compensator for $\multiCountingProc$ if:
 \emph{o.} $\compensator(0)=0$ and $\compensator$ is of finite variation;
 \emph{i.} $\compensator$ is $\filtrationF$-predictable;
 \emph{ii.} $\compensator$ is right-continuous;
 \emph{iii.} $\multiCountingProc - \compensator$ is a local martingale.
\end{defi}

Since a counting process $\multiCountingProc$ is non-decreasing,
as soon as an $\filtrationF$-compensator $\compensator$ for $\multiCountingProc$ exists,
not only is it of finite variation as required by Definition \ref{def.compensator},
but in fact it must be non-decreasing too.
Furthermore, because of requirements \emph{i.} and \emph{ii.},
the discontinuities of $\compensator$ -- if any --
are typically associated with the occurrence of deterministic events,
such as fixed atoms in the distribution of $\multiCountingProc$.
In the following, our concern will be with continuous compensators.

Given the counting process $\multiCountingProc$ and a history $\filtrationF$,
the $\filtrationF$-compensator is unique up to an evanescent set,
and it is equivalently characterised as the $\filtrationF$-predictable projection of
$\multiCountingProc$,
namely as the $\filtrationF$-predictable non-decreasing process $\compensator$ such that
\begin{equation}\label{eq.compensators_and_predict_projections}
 \Expectation \left[ \int_{\R_+} Y d\multiCountingProc\right]
 =
 \Expectation \left[ \int_{\R_+} Y d\compensator\right]
\end{equation}
for all non-negative $\filtrationF$-predictable processes $Y$
(see \cite[Proposition 14.2.II]{DVJ08int}).

If $\compensator$ is absolutely continuous, we write
\begin{equation*}
 \compensator (t) = \intzerot \intensity[ ](s)ds,
\end{equation*}
for some $\filtrationF$-predictable process
$\intensity[ ] = (\intensity[1],\dots,\intensity[\numEventTypes])$,
which is called the intensity of the counting process $\multiCountingProc$.
Predictability is not a technical afterthought:
it is what makes $\intensity(t)$ a function of the strict past,
and every intensity we write down in this chapter will be left-continuous for that reason.
Combining this with equation \eqref{eq.compensators_and_predict_projections},
one obtains the formula
\begin{equation*}
 \Expectation\left[ \countingProc(t)-\countingProc(s) \vert \filtrationF_s \right]
 =
 \Expectation\left[ \int_{s}^{t} \intensity(u) du
  \vert \filtrationF_s \right],
  \qquad  s\leq t,
\end{equation*}
which allows us to interpret $\intensity(t)$ as a measure of the ``instantaneous risk'' of a
jump at time $t$ in the $e$-th component of $\multiCountingProc$.
Notice that this ``risk'' evolves in time,
and that it varies depending on the information available up to time $s$.

We record one consequence separately, because both simulation algorithms of
Section \ref{sec.simulation} rest on it and on nothing else.
Write $\totalIntensity(t) := \intensity[1](t) + \dots + \intensity[\numEventTypes](t)$
for the intensity of the ground process.

\begin{prop}\label{prop.conditionalSurvival}
 Let $\multiCountingProc$ admit the $\filtrationF$-intensity $\intensity[ ]$,
 and let $\arrivalTimes[ ]_{n}$ be the $n$-th arrival of the ground process.
 Then
 \begin{equation}\label{eq.conditionalSurvival}
  \Prob \left( \arrivalTimes[ ]_{n+1} - \arrivalTimes[ ]_{n} > s
  \, \vert \, \filtrationF_{\arrivalTimes[ ]_{n}} \right)
  =
  \exp \left( - \int_{\arrivalTimes[ ]_{n}}^{\arrivalTimes[ ]_{n}+s}
  \totalIntensity(u) \, du \right),
  \qquad s \geq 0,
 \end{equation}
 where $\totalIntensity$ is evaluated on the event that no arrival occurs in the interval.
\end{prop}

On that event the integrand is a deterministic function of the history up to
$\arrivalTimes[ ]_{n}$,
so the right-hand side is $\filtrationF_{\arrivalTimes[ ]_{n}}$-measurable
and \eqref{eq.conditionalSurvival} is a genuine survival function that we can invert.
For a homogeneous Poisson process it reduces to $e^{-\totalIntensity s}$,
and the inter-arrival times are exponential;
in general the whole difficulty of simulating a point process is that
$\int \totalIntensity$ has no closed form,
and the whole point of the kernel we choose in Section \ref{sec.exponentialKernel} is that it does.

Compensators are also the ingredient of the following time-change result,
which we shall use to test whether a simulated path is a path of the model we intended.

\begin{thm}[{\cite{Mey71dem}}]\label{thm.meyer1971}
 Let $\multiCountingProc$ be a $\numEventTypes$-dimensional counting process with arrival times
 $\arrivalTimes$, and assume that
 \emph{i.} no two components of $\multiCountingProc$ jump simultaneously, almost surely;
 \emph{ii.} the compensator $\compensator$ of $\multiCountingProc$ is continuous;
 \emph{iii.} $\compensator_{e}(t) \rightarrow \infty$ as $t\rightarrow \infty$,
 for every $e = 1, \dots, \numEventTypes$.
 Then the time-changed processes
 $\tilde{\countingProc}(u) := \countingProc \big( \inf\lbrace t: \compensator_e(t) > u \rbrace \big)$,
 $e = 1, \dots, \numEventTypes$,
 are mutually independent standard Poisson processes of unit rate.
 Equivalently, the time-changed inter-arrival times
 \begin{equation}\label{eq.time-changed_intertimes}
  \tau_e^j := \compensator(\arrivalTimes_{j}) - \compensator(\arrivalTimes_{j-1})
 \end{equation}
 are independent and exponentially distributed with unit mean,
 jointly over $j = 1, 2, \dots$ and $e = 1,\dots,\numEventTypes$.
\end{thm}

For a proof see \cite{BN88sim}, which also recovers Watanabe's characterisation of the Poisson
process as the special case $\compensator(t) = t$.
Hypothesis \emph{i.} is the one that buys the \emph{joint} independence across components:
the univariate statement, that each $\tilde{\countingProc}$ is separately unit-rate Poisson,
holds without it \citep{Pap72int}.
This distinction is worth keeping in view,
because a goodness-of-fit test that compares each margin against $\text{Exp}(1)$
leaves both of the independences in Theorem \ref{thm.meyer1971} untested
-- the serial independence within a component and the joint independence across components --
and it is exactly a violation of \emph{i.} that such a test cannot see.
